\section{Microscopic variation of the original exceptional union}
\label{sec:v90-fractional}
The new envelope is summable over all microscopic translation scales
at finite collision count. This supplies a genuine regularity statement
on the exceptional source itself. Its topology, and the remaining
pointwise requirement, are specified separately.

For an interval $J$ and $0<s<1$ set
\begin{equation}\label{eq:v90-gagliardo}
 [f]_{s,1;J}=\int_J\int_J
          \frac{|f(t)-f(u)|}{|t-u|^{1+s}}\,du\,dt,\qquad
 \|f\|_{s,1;J}=\|f\|_{L^1(J)}+[f]_{s,1;J}.
\end{equation}
All integrals are actual nonnegative integrals; no sequence of
fixed-scale collision limits defines this seminorm. Put $\chi_0=1/33$.

\begin{theorem}[Complete microscopic oscillation tail]
\label{thm:v90-oscillation-tail}
For $0<s<\chi_0$, fixed $L$ and one fixed auxiliary band,
\begin{align}
 &\sup_{m,R,n,k,a,\delta}\sup_{|J|\le L}m^2
 \int_0^r h^{-1-s}\int_J
 |R_{a,\delta}^\varepsilon(t+h)-R_{a,\delta}^\varepsilon(t)|\,dt\,dh
 \notag\\
 &\hspace{18mm}\le K_{B_*,L}
       \left(\frac{r^{\chi_0-s}}{\chi_0-s}
           +\varepsilon^{-1}\frac{r^{1/2-s}}{1/2-s}\right)
                                                    \label{eq:v90-tail}
\end{align}
for all sufficiently small $r>0$. In particular this complete
microscopic oscillation tail vanishes uniformly before the collision
limsup, with no positive lower translation scale. The corresponding
statement for $b^\varepsilon$ holds for $s<1/32$ with $1/32$ in
place of $\chi_0$.
\end{theorem}
\begin{proof}
Apply \eqref{eq:v90-prefix-envelope} at finite count to a positive
lower truncation of the integral in $h$. Integrate its two explicit
powers. Monotone convergence removes the lower truncation because
$s<\chi_0<1/2$. The same bound holds for every finite count and every
selector, so the supremum may be taken afterwards. Use
\eqref{eq:v90-whole-oscillation} for the last assertion. This order
uses a pre-limit envelope and not an interchange of a collision limsup
with an uncontrolled infinite series.
\end{proof}

\begin{lemma}[Mass and an actual translation envelope]
\label{lem:v90-interpolation}
Suppose $f\in L^1(J)$ and for $0<h<h_0$,
\[
 \int_{J\cap(J-h)}|f(t+h)-f(t)|\,dt\le A h^\chi,
          \qquad 0<s<\chi<1.
\]
If $\mu=\|f\|_{L^1(J)}$, then
\begin{equation}\label{eq:v90-interpolation}
 \|f\|_{s,1;J}\le C_{s,\chi,h_0,|J|}
                 \left(\mu+A^{s/\chi}\mu^{1-s/\chi}\right).
\end{equation}
The convention for $\mu=0$ or $A=0$ is obtained by continuity.
\end{lemma}
\begin{proof}
The integral of a difference is also at most $2\mu$. Changing
variables $u=t+h$ in \eqref{eq:v90-gagliardo} gives twice the
integral of these differences against $h^{-1-s}dh$. Below $h_0$,
use $\min\{2\mu,Ah^\chi\}$ and split at
$(2\mu/A)^{1/\chi}$ when it lies below $h_0$. The two power
integrals are bounded by a constant times
$A^{s/\chi}\mu^{1-s/\chi}$. If the split lies above $h_0$, the
bound is absorbed by a fixed multiple of $\mu$. Above $h_0$ use
$2\mu$. This proves the formula, without an interpolation theorem
for a different norm or a change of source representative.
\end{proof}

Retain exactly the source parameterization used in the revision-89
exceptional-source theorem:
\begin{equation}\label{eq:v90-cofinal-source}
 R_r^\varepsilon=R_{r^{1/8},r^{1/4}}^\varepsilon.
\end{equation}
Its complete interval mass satisfies the inherited ordered estimate
\begin{equation}\label{eq:v90-union-mass-input}
 \limsup_m\sup_{R,n,k}\sup_{|J|\le L}
              m^2\int_JR_r^\varepsilon\le C_L r^{1/128}.
\end{equation}
It follows directly from the critical endpoint source bound and the
capped residence estimate: their leading powers are $r^{129/1024}$
and $r^{1/128}$, respectively. Both sources retain the original
normalization and exact label; their sum is only a positive upper
bound for the union. All auxiliary bands are fixed before the
collision limit and enlarged afterwards.

\begin{theorem}[Fractional regularity and decay on the original union]
\label{thm:v90-original-variation}
For each fixed $\varepsilon>0$, $L>0$ and $0<s<1/33$,
\begin{equation}\label{eq:v90-original-variation}
 \limsup_m\sup_{R,n,k}\sup_{|J|=L}
       \|m^2R_r^\varepsilon\|_{s,1;J}
          \le C_{L,s}(1+\varepsilon^{-1})^{33s}r^{(1-33s)/128}
\end{equation}
for all sufficiently small $r>0$. Moreover the norms
$\|m^2R_{a,\delta}^\varepsilon\|_{s,1;J}$ are bounded at every
finite count by one constant independent of $a,\delta,m,R,n,k$.
These assertions act on the whole original weighted exceptional
union, including its true ends and arbitrary collapsing cutoffs.
\end{theorem}
\begin{proof}
Choose one admissible band once. Equation~\eqref{eq:v90-prefix-envelope}
gives the finite-count translation hypothesis of
Lemma~\ref{lem:v90-interpolation} with $\chi=1/33$ and
$A=K_{B_*,L}(1+\varepsilon^{-1})$, uniformly for all selectors.
The complete unweighted volume bound controls the finite-count
$L^1$ mass since $0\le R_{a,\delta}^\varepsilon\le p$.
This proves the uniform norm assertion.

Apply the lemma at each finite count to the actual function
$m^2R_r^\varepsilon$, and only then take the parameter suprema and
collision limsup. Use \eqref{eq:v90-union-mass-input} for its mass;
that mass estimate may enlarge its own fixed auxiliary band after
the collision limit. The translation constant remains the constant
of the one previously chosen band and has no dependence on $r$.
Since $1-s/\chi=1-33s>0$, the resulting bound is
$C(r^{1/128}+r^{(1-33s)/128})$, as claimed. There is no infinite
scale sum of limiting estimates in this proof: its entire small-scale
part was already controlled by \eqref{eq:v90-tail} at finite count.
\end{proof}

\subsection{High-roof sets and the unchanged essential-height assertion}
For a measurable $E\subset J$, define its regional variational cost
\[
 \operatorname{Cap}_{s,1}(E;J)=
 \inf\{\|u\|_{s,1;J}:u\ge0,\ u\ge1\text{ a.e. on }E\}.
\]
This definition is all that is needed here; no claim about values on
Lebesgue-null sets is intended. Truncation is one-Lipschitz, so taking
$u=\min\{1,m^2R_r^\varepsilon/\lambda\}$ gives, for $\lambda>0$,
\begin{equation}\label{eq:v90-capacity}
 \limsup_m\sup_{R,n,k}\sup_{|J|=L}
 \operatorname{Cap}_{s,1}
       (\{m^2R_r^\varepsilon>\lambda\}\cap J;J)
 \le C_{\varepsilon,L,s}\lambda^{-1}r^{(1-33s)/128}.
\end{equation}
Thus the new estimate measures the spatial oscillation cost of the
large-value set, not just its Lebesgue measure.

It still does not bound the essential height of that set. For example,
on a fixed interval the pulses $\one_{(0,d)}$ have height one,
$L^1$ mass $d$, and fractional seminorm $O(d^{1-s})$ for $s<1$.
They are not a counterexample for the physical source; they explain
why the new regularity topology cannot, by itself, be substituted for
an essential-supremum theorem.

The original scalar target therefore continues to require
\[
 \lim_{B\to\infty}\mathfrak R_{\varepsilon(B),M}=0.
\]
Theorems~\ref{thm:v90-prefix-envelope} and
\ref{thm:v90-oscillation-tail} provide finite-count all-scale
oscillation estimates on exactly that source. Their starting-roof
integral has not been upgraded to a uniform starting-roof bound.
For instance, choosing $r(B)=c\varepsilon(B)^N$ with
$N>4224s/(1-33s)$ makes the explicit bound
\eqref{eq:v90-original-variation} vanish along the legal outer limit.
This is a cofinal fractional-variation conclusion, not charge-zero in
essential height. The prior critical weak, logarithmic, and diffusive-window results
remain proved with their original topologies and normalizations.
No pointwise denominator, same-roof ratio or likelihood is promoted
before scalar essential-height closure and the retained positive
arithmetic hypotheses; zero arithmetic classes are not normalized.
